\subsection{\texorpdfstring{EXPANSIONS OF \(\log(1 + x)\), OF Arc tan x AND OF Arc sin x}{EXPANSIONS OF LOG(1 + X), ARC TAN X AND ARC SIN X}}

Let us integrate the two sides of (26) between 0 and x; we obtain the Taylor expansion of order n of \(\log(1 + x)\) , valid for x \textgreater{} -1

\[
\log (1 + x) = \frac {x}{1} - \frac {x ^ {2}}{2} + \frac {x ^ {3}}{3} + \dots + (- 1) ^ {n - 1} \frac {x ^ {n}}{n} + (- 1) ^ {n} \int_ {0} ^ {x} \frac {t ^ {n} d t}{1 + t}.\tag{27}
\]

The remainder has the same sign as \((-1)^{n}\) if x \textgreater{} 0, and is \textless{} 0 if -1 \textless{} x \textless{} 0; further, when x \textgreater{} 0, we have \(1 + t \geqslant 1\) for \(0 \leqslant t \leqslant x\) , and, when -1 \textless{} x \textless{} 0, we have \(1 + t \geqslant 1 - |x|\) for \(x \leqslant 0\) ; whence the estimates for the remainder

\[
\left| \int_ {0} ^ {x} \frac {t ^ {n} d t}{1 + t} \right| \leqslant \frac {| x | ^ {n + 1}}{n + 1}
\]

\[
\text { for } x \geqslant 0\tag{28}
\]

\[
\left| \int_ {0} ^ {x} \frac {t ^ {n} d t}{1 + t} \right| \leqslant \frac {| x | ^ {n + 1}}{(n + 1) (1 - | x |)} \quad \text {   for   } - 1 <   x \leqslant 0.\tag{29}
\]

From these last two formulae one deduces immediately that for \(-1 < x \leqslant 1\) one has

\[
\log (1 + x) = \sum_ {n = 1} ^ {\infty} (- 1) ^ {n - 1} \frac {x ^ {n}}{n}\tag{30}
\]

the series being uniformly convergent on every compact interval contained in \(]-1,+1]\) , and absolutely convergent for \(|x|<1\) .

On the other hand, for \(|x| > 1\) the general term in the series on the right-hand side of (30) increases indefinitely in size with n (III, p.~106). For x = -1 the series reduces to the harmonic series, which has sum \(+\infty\) (Gen.~Top., IV, p.~365).

Similarly, let us replace \(x\) by \(x^2\) in (26) and integrate both sides between 0 and \(x\) ; we obtain the Taylor expansion of order \(2n - 1\) for \(\operatorname{Arctan} x\) , valid for all real \(x\)

\[
\operatorname{Arc} \tan x = \frac {x}{1} - \frac {x ^ {3}}{3} + \frac {x ^ {5}}{5} + \dots + (- 1) ^ {n - 1} \frac {x ^ {2 n - 1}}{2 n - 1} + (- 1) ^ {n} \int_ {0} ^ {x} \frac {t ^ {2 n} d t}{1 + t ^ {2}}.\tag{31}
\]

The remainder has the sign of \((-1)^n x\) , and since \(1 + t^2 \geqslant 1\) for all \(t\) we have the estimate

\[
\left| \int_ {0} ^ {x} \frac {t ^ {2 n} d t}{1 + t ^ {2}} \right| \leqslant \frac {| x | ^ {2 n + 1}}{2 n + 1}\tag{32}
\]

from which one deduces that, for \(|x| \leqslant 1\) ,

\[
\operatorname{Arc} \tan x = \sum_ {n = 1} ^ {\infty} (- 1) ^ {n - 1} \frac {x ^ {2 n - 1}}{2 n - 1}\tag{33}
\]

the series being uniformly convergent on \([-1,+1]\) , and absolutely convergent for \(|x|<1\) .

In particular, for \(x = 1\) one obtains the formula

\[
{\frac {\pi}{4}} = 1 - {\frac {1}{3}} + {\frac {1}{5}} + \dots + (- 1) ^ {n} {\frac {1}{2 n + 1}} + \dots .\tag{34}
\]

For \(|x| > 1\) the general term of the right-hand side of (33) increases indefinitely in size with \(n\) .

Finally, for the Taylor expansion of Arc sin x we start from the expansion of its derivative \((1 - x^{2})^{-1/2}\) ; this last expansion is obtained by replacing x by \(-x^{2}\) in the expansion of \((1 + x)^{-1/2}\) as a binomial series; for \(|x| < 1\) this gives

\[
(1 - x ^ {2}) ^ {- 1 / 2} = 1 + \frac {1}{2} x ^ {2} + \frac {1 . 3}{2 . 4} x ^ {4} + \dots + \frac {1 . 3 . 5 \dots (2 n - 1)}{2 . 4 . 6 \dots 2 n} x ^ {2 n} + r _ {n} (x)
\]

with, by (23), the bound

\[
0 \leqslant r _ {n} (x) \leqslant \frac {x ^ {2 n + 2}}{\sqrt {1 - x ^ {2}}}
\]

for the remainder.

On taking the primitive of the preceding expansion we obtain

\[
\operatorname{Arc} \sin x = x + \frac {1}{2} \frac {x ^ {3}}{3} + \frac {1 . 3}{2 . 4} \frac {x ^ {5}}{5} + \dots + \frac {1 . 3 . 5 \dots (2 n - 1)}{2 . 4 . 6 \dots 2 n} \frac {x ^ {2 n + 1}}{2 n + 1} + \mathrm{R} _ {n} (x)\tag{35}
\]

where \(\mathrm{R}_{n}(x)\) has the sign of x and satisfies the inequality

\[
| \mathrm{R} _ {n} (x) | \leqslant \int_ {0} ^ {x} \frac {t ^ {2 n + 2} d t}{\sqrt {1 - t ^ {2}}}.\tag{36}
\]

Further, the relation (35) shows that \(\mathbf{R}_n(x)\) tends to a limit when \(x\) approaches 1 or \(-1\) , so one has

\[
| \mathrm{R} _ {n} (1) | \leqslant \int_ {0} ^ {1} \frac {t ^ {2 n + 2} d t}{\sqrt {1 - t ^ {2}}}.\tag{37}
\]

But the right-hand side of (37) tends to 0 when \(n\) tends to \(+\infty\) : for, since the integral \(\int_0^1 dt / \sqrt{1 - t^2}\) is convergent, for every \(\varepsilon > 0\) there is an \(a\) such that \(0 < a < 1\) and \(\int_a^1 dt / \sqrt{1 - t^2} \leqslant \varepsilon\) ; on the other hand we have

\[
\int_ {0} ^ {a} \frac {t ^ {2 n + 2} d t}{\sqrt {1 - t ^ {2}}} \leqslant \frac {1}{\sqrt {1 - a ^ {2}}} \int_ {0} ^ {a} t ^ {2 n + 2} d t = \frac {a ^ {2 n + 3}}{(2 n + 3) \sqrt {1 - a ^ {2}}}
\]

and so there is an \(n_0\) such that for \(n \geqslant n_0\) one has \(\frac{a^{2n+3}}{(2n+3)\sqrt{1-a^2}} \leqslant \varepsilon\) , whence, finally, \(|\mathrm{R}_n(x)| \leqslant 2\varepsilon\) for \(|x| \leqslant 1\) and \(n \geqslant n_0\) . Thus one has

\[
\operatorname{Arc} \sin x = \sum_ {n = 0} ^ {\infty} \frac {1 . 3 . 5 \dots (2 n - 1)}{2 . 4 . 6 \dots 2 n} \frac {x ^ {2 n + 1}}{2 n + 1}\tag{38}
\]

the right-hand side being normally convergent on the compact interval \([-1, 1]\) .

In the opposite case one can show, as for the binomial series, that the general term in the series on the right-hand side of (38) increases indefinitely in absolute value for \(|x| > 1\) .

On putting \(x = \frac{1}{2}\) , for example, in (38) we obtain a new expression for the number \(\pi\) ;

\[
\frac {\pi}{6} = \sum_ {n = 0} ^ {\infty} \frac {1 . 3 . 5 \dots (2 n - 1)}{2 . 4 . 6 \dots 2 n} \frac {1}{(2 n + 1) 2 ^ {2 n + 1}}
\]

which is much better suited than formula (34) to calculating approximations to \(\pi\) (see Calcul numérique); one thus obtains

\[
\pi = 3. 1 4 1   5 9 2   6 5 3 \dots
\]

accurate to within \(1/10^{9}\) . \(^{3}\)
% semantic-fragments: legacy-input-seam
\relax{}
\setcounter{exercisegroup}{0}
\refstepcounter{exercisegroup}
